% ----------------------------------------------------------------------
% Why:     Android teams increasingly want the app engineer to own the API contract too.
%          This variant leads with the largest Android codebase and then shows the
%          backend work as backend work, without dressing it up as mobile.
% What:    One-page resume targeting Android roles, with the backend APIs behind
%          the app as supporting evidence.
%          Built from the resume-forge evidence ledger; every number traces to a
%          repository.
% Result:  A 34-module Compose codebase with 199 test files first, and honest
%          provenance on the training project that follows it.
% Changed: 2026-09-27 -- created from assets/resume-template.tex.
% Build:   bash ~/.claude/skills/resume-forge/scripts/build_resume.sh \
%            public/resume-mobile-developer.tex
% Check:   python3 ~/.claude/skills/resume-forge/scripts/evaluate_resume.py \
%            public/resume-mobile-developer.tex --role mobile-developer --min-score 85
% ----------------------------------------------------------------------
\documentclass[10pt,a4paper]{article}

% ----------------------------------------------------------------------
% Packages
% ----------------------------------------------------------------------
\usepackage[T1]{fontenc}
\usepackage[utf8]{inputenc}
\usepackage{lmodern}
\usepackage[a4paper,top=0.9cm,bottom=0.8cm,left=1.25cm,right=1.25cm]{geometry}
\usepackage{titlesec}
\usepackage{enumitem}
\usepackage[hidelinks]{hyperref}

\hypersetup{
  pdftitle={Harish G - Android Developer Resume},
  pdfauthor={Harish G},
  colorlinks=false
}

% ----------------------------------------------------------------------
% Section styling: sans-serif bold heading with a thin rule beneath
% ----------------------------------------------------------------------
\titleformat{\section}
  {\normalsize\sffamily\bfseries}{}{0em}{}[\titlerule]
\titlespacing*{\section}{0pt}{8pt}{3pt}

% Compact bullet lists: the single biggest lever on whether this fits a page.
\setlist[itemize]{leftmargin=1.2em,itemsep=1.5pt,topsep=2pt,parsep=0pt}

\setlength{\parindent}{0pt}
\linespread{0.97}\selectfont
\pagestyle{empty}

% ----------------------------------------------------------------------
% Helpers
% ----------------------------------------------------------------------
% Entry heading row with a right-aligned date: \entry{Title}{Date}
\newcommand{\entry}[2]{\textbf{#1}\hfill{#2}\par}
% Middot separator: keeps the skills block scannable instead of prose-like.
\newcommand{\sep}{\,\textperiodcentered\,}
% Fixed-width label so every skills row starts at the same column. A fixed
% \makebox overlapped its own content whenever a label was wider than the box,
% so the gap is measured instead: aligned when the label fits, and a plain 0.6em
% gap when it does not. A long label can never collide with the row again.
\newlength{\skillwidth}\setlength{\skillwidth}{6.4em}
\newlength{\skilllabel}
\newcommand{\skill}[2]{%
  \settowidth{\skilllabel}{\textbf{#1}}%
  \hangindent=\skillwidth\hangafter=1%
  \mbox{\textbf{#1}}%  \mbox so inner spaces cannot stretch under justification
  \ifdim\skilllabel<\skillwidth\hspace{\dimexpr\skillwidth-\skilllabel}\else\hspace{0.6em}\fi
  #2\par}

\begin{document}

% ----------------------------------------------------------------------
% Header
% ----------------------------------------------------------------------
{\LARGE\bfseries Harish G}\\[1pt]
{\normalsize\bfseries Android Developer \textbar{} Kotlin, Jetpack Compose, Room, Hilt \textbar{} and the APIs behind the app}\\[2pt]
{\small Bengaluru, India \quad\textbar\quad +91 7892855850 \quad\textbar\quad
\href{mailto:harishgreddy.work@gmail.com}{harishgreddy.work@gmail.com} \quad\textbar\quad
\href{https://www.linkedin.com/in/harishgreddy/}{LinkedIn} \quad\textbar\quad
\href{https://github.com/harish00506/}{GitHub}}

\section{Profile}
I build Android apps and the services behind them. That makes me the one who notices when the API
contract and the screen disagree. My largest Compose codebase runs to 34 Gradle modules and 199 test
files, and the calculation engines are pure Kotlin, so they test without an emulator. The backend
half is streaming. It is the same reconnect problem a phone on a bad network has.

\section{Core Skills}
\skill{Mobile}{Kotlin \sep Jetpack Compose \sep MVVM \sep Hilt \sep Room / SQLCipher \sep ML Kit OCR \sep Firebase Auth \& Realtime DB \sep Gradle Modules \sep Android Studio}
\skill{Languages}{Kotlin \sep Java \sep Python \sep TypeScript \sep JavaScript \sep SQL}
\skill{Backend}{FastAPI \sep Spring Boot \sep Node.js/Express \sep REST API Design \sep WebSockets \sep Celery}
\skill{Data}{Room / SQLite \sep Firebase \sep PostgreSQL \sep MongoDB \sep Redis}
\skill{Testing}{JVM unit tests \sep pytest \sep GitHub Actions \sep GitLab CI \sep Docker \& Compose}
\skill{Concepts}{Modular Architecture \sep One-Way Module Graph \sep On-Device Processing \sep JWT Auth \& RBAC \sep System Design}

\section{Experience}
\entry{Software Developer, CortexCraft.ai}{Jan 2026 -- Present}
\begin{itemize}
  \item Built and maintain the voice-agent platform (589 of 626 commits on the SaaS platform, 67 of 75
        on the core engine): real-time audio streaming over WebSockets, and an external REST API that
        client applications drive calls through.
  \item Own production reliability of that pipeline: fixed outbound routing, added reconnect handling
        for dropped sockets, and traced a Docker TLS failure to NAT hairpinning. Reconnect behaviour
        on a flaky network is the same problem a mobile client has.
  \item Shipped LeadCall AI on the platform: multi-tenant outbound calling with Celery workers,
        Alembic-migrated PostgreSQL and a React dashboard.
\end{itemize}

\vspace{1pt}
\entry{Freelance Web \& App Developer (Self-Employed)}{2023 -- 2026}
\begin{itemize}
  \item Delivered web and app projects end to end for local clients, including Android work, React front ends and the services behind them.
\end{itemize}

\section{Selected Projects}
\textbf{AI Personal CFO and Talo} \textit{Kotlin, Jetpack Compose, Hilt, Room/SQLCipher, ML Kit, Gemini Live}
\begin{itemize}
  \item Sole author, 159 commits, 34 Gradle modules, 199 test files. A one-way module graph keeps
        feature modules independent of each other, and the calculation engines are pure Kotlin with
        no Android imports, so they are testable without an emulator.
  \item Compose UI over immutable state, Hilt constructor injection throughout, and Room on SQLCipher
        with the database key held in the Android Keystore. On-device ML Kit OCR extracts receipt
        fields, and the SMS parser is built to refuse, resolving every ambiguous case to null.
  \item Twelve deterministic engines compute every figure and the model is only permitted to explain
        them; an output guardrail rewrites any number absent from the tool result as
        \texttt{[unverified]}.
\end{itemize}

\vspace{1pt}
\textbf{GroceryGo} \textit{Kotlin, Jetpack Compose, Firebase}
\begin{itemize}
  \item 4 of 4 commits. Android grocery-list app with a Jetpack Compose UI and ViewModel-held screen
        state on an MVVM structure, backed by Firebase Authentication and Realtime Database.
  \item Final project of my Internshala Android development training, and where I learned MVVM on
        Android.
\end{itemize}

\vspace{1pt}
\textbf{KisanVoice AI} \textit{Python, FastAPI, React, MongoDB, WhatsApp Flows, Google STT/TTS}
\begin{itemize}
  \item 172 of 192 commits. Mobile-first survey platform reaching users through WhatsApp rather than
        an installed app: voice answers in 24 configured locales (10 Indian), with WhatsApp Flows
        standing in for long option sets that would be unusable as a text menu.
\end{itemize}

\section{Education \& Ways of Working}
\textbf{B.E., Information Science and Engineering}, Vivekananda Institute of Technology, Bangalore\\
Feature branches and conventional commits; GitLab CI with 132 pytest files as the merge gate on the
voice-agent platform.

\end{document}
